\section*{Part II: The Anatomy of Cost}

\subsection*{6. The Four Costs a Founder Must Separate}

\textbf{Fixed costs within a relevant range} do not change materially with 1 additional customer during the planning period: creating a curriculum, maintaining a base platform, preparing a cohort, a committed staff salary, or minimum service infrastructure. Fixed does not mean permanent, optional, or irrelevant to long-run profit.

\textbf{Variable costs} change with usage or customers: payment processing, learner-specific supplies, cloud inference, individually delivered feedback, shipping, or a per-customer software license. The relationship need not be linear.

\textbf{Step-fixed costs} remain stable until a threshold is crossed, then jump: hiring another reviewer, renting an additional classroom, adding a 2nd support shift, or purchasing a 2nd server. The next learner may cost almost nothing in 1 range and require an entire new capacity block in the next.

\textbf{Committed or sunk costs} have already been incurred or cannot be recovered economically at the decision moment. They matter when evaluating past performance, but generally should not determine whether to accept an incremental profitable order. Future avoidable and opportunity costs matter more for that decision.

The classification is always \textbf{relative to a decision, time horizon, and volume range}. An instructor's salary is fixed over the current month, but may be variable when planning next year's staffing. A software subscription may be fixed until another seat is needed. ``Fixed'' is not an eternal property of an expense.

\subsection*{7. The Cost Function and Why It Is Often Not a Straight Line}

For a simple product, write:

\begin{equation}
C(Q) = F + cQ,
\end{equation}
where $Q$ is units delivered, $F$ is fixed cost during the relevant range, and $c$ is constant incremental cost. Average cost is $F/Q+c$; fixed cost is spread as quantity rises.

Real services often require a more honest representation:

\begin{equation}
C(Q) = F + cQ + \sum_j K_j\,\mathbf{1}\{Q>q_j\} + R(Q) + V(Q),
\end{equation}
where $K_j$ is a discrete capacity addition after threshold $q_j$, $R(Q)$ denotes quality-related rework and exceptions, and $V(Q)$ other nonlinear effects such as demand variability, specialist escalation, or overtime. These are planning terms, not quantities 1 can generally estimate precisely before operating.

The cost curve can flatten through reuse, jump with hiring, and rise sharply when queues and errors appear. That is why ``our costs remain the same if we double the cohort'' is an unusually strong claim. It requires evidence that the original service standard survives the increase.

\subsection*{8. Marginal Cost Is a Decision Concept, Not a Slogan About Software}

Marginal cost is the additional cost of producing 1 more unit, evaluated under relevant constraints:

\begin{equation}
MC(Q) \approx C(Q+1) - C(Q).
\end{equation}

For a downloadable video, the marginal distribution cost may be very low. For admission to a supervised course, the marginal \emph{complete-service} cost includes expected reviewing, support, infrastructure, payment, and failure-recovery work. If the next customer pushes the firm past a capacity threshold, the immediate incremental cost may be extremely high unless unused capacity from the addition will be sold later.

This is why \textbf{digital content has 1 marginal-cost curve, personalized learning has another, and the combined product has a 3rd}. Low-cost copies of content do not eliminate the costly work that gives the content commercial value to some buyers.

For an incremental sale, compare incremental contribution with incremental costs and foreseeable effects on existing customers. Do not allocate every historical fixed expense to a single additional customer and declare the sale unattractive. But do not assume that existing instructor capacity is free if it is already occupied or needed to protect service quality.

\subsection*{9. Contribution, Operating Profit, and Cash Are Different Views}

For a simplified offering with realized net fee $P$, variable acquisition and fulfillment cost $v$ per customer, and fixed launch-and-delivery cost $F$:

\begin{equation}
\Pi_{\text{cohort}} = N(P-v) - F,
\end{equation}
where $\Pi_{\text{cohort}}$ is modeled cohort contribution after dedicated fixed delivery costs.

Some firms would call $N(P-v)$ contribution margin and subtract fixed costs below it; naming conventions vary. Specify which costs have been deducted rather than relying on a label.

This expression is \textbf{not} automatically accounting profit, economic profit, or cash generated. The company may owe taxes, have other overhead, capitalize certain investments under its accounting policy, hold prepaid tuition against future services, or wait months to collect a B2B invoice. Nor is it sustainable if the founder performs valuable labor without being counted.

A 1st-stage entrepreneur should keep parallel views: \textbf{cash economics for survival; normalized economics for repeatability; capacity economics for reliability; and outcome economics for value}. Each can contradict the others without being false.

\subsection*{10. The Price Waterfall Begins Before the First Lesson}

A tuition fee of ₹30,000 does not guarantee that ₹30,000 becomes retained revenue. Enrollment discounts, approved refunds, tax treatment, transaction fees, bad debts, and promotional concessions affect the realized amount. Acquisition effort may be substantial. Even direct inbound inquiries require the founder's time.

A useful operating model separates list price, realized fee excluding taxes held for authorities, gross collections, expected refunds and defaults, per-sale transaction costs, acquisition cost, variable delivery cost, and cohort overhead. Treat outstanding delivery obligations separately from the money in the bank.

This chapter is about fulfillment, but the economics of fulfilling a service cannot be understood without its acquisition and payment conditions. A product with excellent classroom costs and ruinous sales costs is not a healthy business.
